\section{A first-root carrier: fifteen integers, degree 56}
\label{sec:first-carrier}
The principal new equation has \textbf{15 unrestricted integer unknowns,
exact degree 56, and 110 live binary arithmetic operations}.  It improves
both degree and evaluator cost over the inspected 15-variable degree-170
and degree-172 alternatives.  The construction is not a generic variable
renaming: it makes the first Pell root carry two coprime congruences, replacing
the separate main-root coordinate and a large radical norm.

The unknowns are
\[
 j,C,z,h,r,v,w,T,c,n,f,y,p_1,p_2,p_3\in\mathbb Z.
\]
Use the same six fixed source34 coefficients as in the retained calibrated
source, with $M_F$ already shifted.  All expressions below are abbreviations,
not additional unknowns:
\begin{align*}
 J&=6j+1,& q&=\beta J+1,& W&=12w+4,\\
 X&=qW,&Y&=qv,&F&=(K_0+W)C-z(q-1),\\
 R&=(q(q-F)-C)(q^2-1)+(M_C+qM_F)J,\\
 E&=XY,&a&=E+Y,&H&=4a+3,\\
 \Delta&=a^2+H,&k&=R+1+hE,&\eta&=c-Yk,\\
 L&=EYk,& I&=(2r+1)c^2,&Q&=\Delta(f^2-1),\\
 V&=f^2(cT-R)-c,& M&=2q^2,& U&=MH,\\
 K_*&=2X^2-3Xc-2,&Z&=Un+MK_*+H(v-q^2).
\end{align*}
Put
\begin{align*}
 D&=(\mathrm{cone}_A F-q)^2+(\mathrm{cone}_B C-q)^2,\\
 B_{\rm sign}&=k^2(q^2-D)-4q^2\eta^2,\\
 P_{\rm pos}&=1+p_1^2+p_2^2+p_3(p_3+1),\\
 A_1&=Z^2+(2Z-Uk)UL-U^2,\\
 A_2&=f^2-\Delta I^2-1,\\
 A_3&=V^2-(Q-1)y(2V+y)-1,\\
 A_4&=B_{\rm sign}-P_{\rm pos}.
\end{align*}
The complete equation is
\begin{equation}\label{eq:carrier15}
 \boxed{P_{15}=A_1^2+A_2^2+A_3^2+A_4^2=0.}
\end{equation}
Its literal DAG is \texttt{lowdegree/equation.json}.

\paragraph{Why the compression is sound.}
Unconditionally, $q$ is even, $q\equiv2\pmod3$, $H$ is odd, and
$H\equiv3\pmod q$.  Hence $\gcd(H,M)=1$ and $U\ne0$.
The first row makes $t=Z/U$ a rational root of the monic integer polynomial
$t^2+(2t-k)L-1$, so $t\in\mathbb Z$.  Coprime CRT gives
\[
 v=q^2(2v_0+1),\qquad H\mid K_*.
\]
Thus the cubic scale is recovered \emph{before} the signed Pell bootstrap.
For $K=X(X+2ac)-1$, the identity $2K-K_*=HXc$ implies $H\mid K$.
The budget row supplies the common native signs and $|\eta/k|<1/2$.
The strong and auxiliary rows recover the omitted main Pell square by the
ordinary folded-index lemma.  Normalized rank then gives its index $R$.
The retained whole-product estimate
$|H|>2^{2R+2}|X|^2$ turns $H\mid K$ into the exact signed power relation,
and restores an integral main coordinate
$\gamma=(d-X-ac)/H$.  This yields a literal zero of the pinned normalized
17-variable source.  The proof does not assume positivity of $\eta$ before
the source's centered half-binomial argument recovers it.

Conversely, every zero of that source has even $q$, hence even old $J_3$.
Set $j=J_3/2$, $C=q-F-\alpha$, $v=q^2(2v_0+1)$, and
$n=t_{\rm old}-K_*/H-v_0$.  The old main row proves $H\mid K_*$;
all the displayed coordinates are integers and all four new rows vanish.
The complete signed proof and maps are in \texttt{lowdegree/PROOF.md}.

\paragraph{Exact, not merely syntactic, size.}
The row degrees are $28,14,18,16$.  The actual fixed-coefficient polynomial has
\[
 [j^{30}w^{10}v^{12}h^4]P_{15}=2^{62}3^{40}\beta^{30}>0.
\]
Thus its exact degree is56.  Its110 live gates are63 multiplications,
24 additions and23 subtractions.  Independent normal and optimized-Python
checks verify all register cuts, coefficients, liveness, the leading form,
and rejection of corrupted gates.

\section{A double-centered companion: fourteen integers, degree 108}
\label{sec:double-carrier}
The second equation removes the explicit $f$, leaving \textbf{14 unrestricted
integers, exact degree108, and119 live operations}.  Keep the preceding
abbreviations except replace $f^2$ everywhere by
\[
 D_f=1+\Delta I^2,\qquad Q=\Delta(D_f-1),\qquad
 V=D_f(cT-R)-c.
\]
Define the two doubled root centers and errors
\begin{align*}
 A&=2L+k,&B&=2(a+2)I,\\
 a_*&=4-k^2,&b_*&=4(1-I^2),\\
 D_0&=Z^2+2BUZ+U^2(a_*-b_*),\\
 E_0&=Z+U(A+B),\\
 B_0&=A[Z^2+2BUZ-b_*U^2]-a_*U[Z+UB],\\
 N&=D_0^2+4UE_0B_0.
\end{align*}
Now $A^2+a_*=4D_\tau$, $B^2+b_*=4D_f$, where
$D_\tau=1+L(L+k)$.  The complete equation is
\begin{equation}\label{eq:carrier14}
 \boxed{P_{14}=N^2+[V^2-(Q-1)y(2V+y)-1]^2+
 [B_{\rm sign}-P_{\rm pos}]^2=0.}
\end{equation}
The exact two-root norm identity makes $Z/U+A+B$ a rational root of a monic
quartic over $\mathbb Z$.  Hence CRT again recovers the scale first, without
assuming radical rationality or excluding $v=0$ by hand.  The guard and
recovered scale give $\sqrt{D_f}>2\sqrt{D_\tau}$, so the nonzero rational
signed root sum forces both roots rational and therefore integral.  This
restores the15-variable equation by selecting $f=\sqrt{D_f}$ and an integral
new $n$.  Conversely every15-variable zero gives a zero here by shifting
$n$ to the sum of its doubled integer roots minus $A+B$ and the CRT offsets.

The fully centered norm has degree54, not the larger bound obtained by
expanding an uncentered norm syntactically.  In the final polynomial,
\[
 [j^{56}w^{18}v^{22}h^4r^2c^4n^2]P_{14}
   =2^{130}3^{74}\beta^{56}>0.
\]
Its119 live gates are69 multiplications,30 additions and20 subtractions.
The literal equation, full proof, and separate independent review are under
\texttt{lowdegree/fourteen/}.

\paragraph{Common semantic boundary.}
Both reductions use the unchanged authenticated source34 coefficient tuple
and the retained ordinary normalized-source, native compiler and finite-ZFC-
refutation calibration.  They do not prove consistency, claim global
minimality, or supply a new end-to-end proof-assistant or encoded internal-ZFC
certificate.  Their two-way arithmetic maps are explicit, with no hidden
existential registers or domain restrictions.
